% Source: https://artofproblemsolving.com/community/c3h1605793_2018_mathcounts_school_sprint_12
% Stable identity: aops-topic-1605793-q01
Let the square have side $s$; the original fence has length $4s$. The new fence joins the midpoints of two opposite sides, so it is parallel to the other two sides and has length $s$:

\begin{asy}
usepackage("amsmath");
size(4cm);
defaultpen(fontsize(10pt));
real s=4;
draw((0,0)--(s,0)--(s,s)--(0,s)--cycle,linewidth(1));
draw((s/2,0)--(s/2,s),red+linewidth(1)+dashed);
dot((s/2,0),red);dot((s/2,s),red);
label("$s$",(s/2,0),S);label("$s$",(0,s/2),W);label("$s$",(s/2,s),N);label("$s$",(s,s/2),E);
label("new fence",(s/2,s/2),E,red);
label("\mbox{}",(min(currentpicture,user=true)+max(currentpicture,user=true))/2,black); // keeps the text after the figure black when LaTeX typesets the labels inline
\end{asy}

The fencing grows from $4s$ to $5s$, an increase of $\dfrac{s}{4s}=\dfrac{1}{4}$, that is, $\boxed{25\%}$.
